

\vspace{-1mm}
\section{Preliminaries}
\vspace{-2.5mm}

% In this section, we first present the fundamentals of RLVR. We then introduce the pass@$k$ metric, which we use to evaluate the reasoning boundary. We also explain why we choose pass@$k$ over alternative evaluation methods such as best-of-$N$ or majority voting.

In this section, we first outline the fundamentals of RLVR, then introduce the pass@$k$ metric to evaluate reasoning boundaries, and explain why it is preferred over alternatives like best-of-$N$.



\begin{figure}[t!]
% \vspace{-2mm}
  \centering
   % \includegraphics[width=0.95\textwidth]{tex/images/simpleRL_coverage_verynew.pdf}
   % \includegraphics[width=0.95\textwidth]{tex/images/simpleRL_coverage_verynew_v2.pdf}
      \includegraphics[width=0.9\textwidth]{tex/images/simpleRL_coverage_main.pdf}
\captionsetup{font={footnotesize}}
\caption{
Pass@$k$ curves of base models and their RLVR-trained counterparts across multiple mathematical benchmarks.  
When $k$ is small, RL-trained models outperform their base versions.  
However, as $k$ increases to the tens or hundreds, base models consistently catch up and surpass RL-trained models. More results on GSM8K and AMC23 can be found at~\cref{fig:mat_appendix}.
}
\label{fig:math}
\vspace{-5mm}
\end{figure}


\vspace{-1.5mm}
\subsection{Reinforcement Learning with Verifiable Rewards}
\vspace{-1.5mm}
\label{sec:prelim}


\textbf{Verifiable Rewards.} Let $\pi_{\theta}$ be an LLM with parameters~$\theta$ that generates a token sequence
$\mathbf{y} = (y_{1}, \dots, y_{T})$ conditioned on a natural-language prompt~$x$.
A deterministic \emph{verifier} $\mathcal{V}$ returns a binary reward: $r = \mathcal{V}(x, \mathbf{y}) \in \{0, 1\},$
where $r = 1$ if and only if the model's final answer is exactly correct.
% (\eg, the numeric result matches the ground truth in math, or all unit tests pass in code)
% .
A format reward may also be added to encourage the model to explicitly separate the reasoning process from the final answer.
% Each prompt-response pair is treated as a single-step episodic Markov Decision Process (MDP). 
The goal of RL is to learn a policy to maximize the expected reward:
$
J(\theta) =
\mathbb{E}_{x \sim \mathcal{D}}
\left[
\mathbb{E}_{\mathbf{y} \sim \pi_{\theta}(\cdot | x)} [\, r \,]
\right]
$
, where $\mathcal{D}$ is the distribution of prompts.

\textbf{RLVR Algorithms.} 
Proximal Policy Optimization (PPO)~\cite{schulman2017ppo} proposed using the following clipped surrogate to maximize the objective:
\begin{equation}
\setlength\abovedisplayskip{2pt}
\setlength\belowdisplayskip{2pt}
\mathcal{L}_{\text{CLIP}} =
\mathbb{E} \left[
\min\!\left(
r_t(\theta) A_t,\;
\operatorname{clip}(r_t(\theta), 1 - \epsilon, 1 + \epsilon) A_t
\right)
\right],
\end{equation}
where $r_t(\theta) =
\pi_{\theta}(y_t | x, \mathbf{y}_{<t}) /
\pi_{\theta_{\text{old}}}(y_t |  x, \mathbf{y}_{<t})$,
and $A_t$ is the advantage estimated by a value network~$V_{\phi}$. 
KL divergence term is optionally applied, 
% either as a reward penalty or directly in the loss function, 
to constrain the model from deviating too far from the original policy.
More algorithms are introduced in~\cref{appendix:algos}.

\textbf{Policy Gradient.}  
PPO and its variants belong to the policy gradient class of RL ~\cite{williams1992reinforce,sutton1998reinforcement}. These methods learn exclusively from \textit{on-policy samples}, \ie, samples generated by the current LLM.  
In the context of verifiable rewards, the training objective generally \textit{maximizes the log-likelihood of samples with correct answers and minimizes the likelihood of those with incorrect answers.}


\textbf{Zero RL Training} applies RL directly to the base model without any supervised fine-tuning (SFT)~\cite{guo2025deepseek-r1}.
To clearly study the effect of RLVR, we follow this zero-RL setting for all math tasks using pretrained models as start model.
However, for coding and visual reasoning tasks, open-source work typically uses instruction-tuned models as starting points, primarily due to the training instability and limited effectiveness of using a pure zero-RL setting.
Following this convention, we compare the finetuned model with its RLVR-trained counterpart to focus solely on the effect of RLVR.
% We follow this convention, which is reasonable for studying the effect of RL on reasoning ability, as these models are primarily fine-tuned on general instruction-following data rather than long-form CoT reasoning data.

\vspace{-1.5mm}
\subsection{Metrics for LLM Reasoning Capacity Boundary}
\vspace{-1.5mm}

\textbf{Pass@$k$ Metrics.}
Accurately measuring the reasoning ability boundary of base and RL models is challenging, as methods like greedy decoding or the average of nucleus samplings~\cite{holtzman2019topp} only reflect average-case performance.
% For example, consider a traditional avg@8 approach, where 8 responses are sampled for an extremely difficult problem. If none of the responses are correct, the value is 0/8. However, after more trials such as 128, the model may solve the problem, indicating that the reasoning boundary can eventually overcome this challenge.
To accurately measure the reasoning ability boundary, 
we extend the commonly used pass@\textit{k} metric from code generation~\cite{chen2021evaluating} to all tasks with verifiable rewards. Given a problem, we sample $k$ outputs from the model. 
The pass@$k$ value for this question is 1 if at least one of the $k$ samples passes verification; otherwise, it is 0.
The average pass@$k$ value over the dataset reflects the proportion of problems in the dataset that the model can solve within $k$ trials, providing a rigorous evaluation of the reasoning capacity coverage of LLMs.
We adopt an unbiased, low-variance estimator for computing to calculate pass@$k$, as detailed in~\cref{appendix:passk_est}.

\textbf{Comparison with Best-of-$N$ and Majority Voting.}
Best-of-$N$~\cite{cobbe2021gsm8k} and majority voting are practical methods for selecting correct answers, but they may overlook a model's full reasoning potential. In contrast, we use pass@$k$ \textit{not to assess practical utility} but to investigate the boundaries of reasoning capacity. 
If a model produces a correct solution in any of the $k$ samples, we treat the problem as within its potential scope. Thus, if RL enhances reasoning, the RL-trained model should succeed in more such problems than the base model. 
Methods like Best-of-$N$ or majority voting may miss these successes if the correct answer is not selected by the verifier or voting.


\textbf{Random Guessing Issue}. 
For \textit{coding} tasks, where a compiler and predefined unit test cases are used as verifiers, the pass@$k$ value can accurately reflect whether the model can solve the problem. 
In \textit{mathematics}, the issue of ``guessing'' can become pronounced as $k$ increases, where a model may generate an incorrect CoT but still accidentally arrive at the correct answer.
To address this, we manually check the correctness of CoT for a subset of model outputs as detailed in~\cref{sec:math}. 
By combining results on math with manually checking and coding, we rigorously evaluate the scope of LLM's reasoning capacity. 
Another caveat is that, with an astronomically large $k$, even uniform sampling over the token dictionary would stumble upon the correct reasoning path--though this is infeasible within today's time and compute budgets. Crucially, we find that the base model already produces correct outputs at realistic values ($k = 128$ or $1024$), well within practical resource limits.



\begin{table*}[ht]
\centering
\small
% \scriptsize
\renewcommand{\arraystretch}{1.15}
\setlength{\tabcolsep}{3pt}
\caption{Experimental setup for assessing RLVR's effect on the reasoning boundaries of LLMs.}
\label{tab:exp-settings}

\resizebox{0.98\textwidth}{!}{
    % \begin{tabularx}{\linewidth}{@{}lX X X X@{}}
    \begin{tabular}{ccccc}
    % \toprule
    \textbf{Task} & \textbf{Start Model} & \textbf{RL Framework} & \textbf{RL Algorithm(s)} & \textbf{Benchmark(s)} \\ %\vspace{-0.8mm}
    \midrule

    % \noalign{\vskip -1mm}
    \textbf{Mathematics} & 
    \makecell[c]{LLaMA-3.1-8B \\ Qwen2.5-7B/14B/32B-Base \\ Qwen2.5-Math-7B }& 
    \makecell[c]{SimpleRLZoo \\ Oat-Zero \\ DAPO}  & 
    GRPO & 
    \makecell[c]{GSM8K, MATH500\\ Minerva, Olympiad \\ AIME24, AMC23}  \\

    % \midrule
    \rowcolor{gray!15}\textbf{Code Generation} & 
    \makecell[c]{ Qwen2.5-7B-Instruct \\ DeepSeek-R1-Distill-Qwen-14B}  & 
   \makecell[c]{ \ \ \ \ \ \ \ \ \ Code-R1 \ \ \ \ \ \ \ \ \ \ \\ DeepCoder} & 
    GRPO & 
    \makecell[c]{\ \ \ \ \ LiveCodeBench\ \ \ \ \  \\ HumanEval+ }
   \\

    % \midrule
    \textbf{Visual Reasoning} & 
    Qwen2.5-VL-7B & 
    EasyR1 & 
    GRPO & 
    \makecell[c]{MathVista \\ MathVision} \\
    

    % \midrule
    \rowcolor{gray!15}\textbf{Deep Analysis} & 
    \makecell[c]{Qwen2.5-7B-Base  \\ Qwen2.5-7B-Instruct \\ DeepSeek-R1-Distill-Qwen-7B \  \ } & 
     VeRL & 
    \makecell{PPO, GRPO \\ Reinforce++ \\ RLOO, ReMax, DAPO}  & 
    \makecell[c]{\ \ \ \ Omni-Math-Rule\ \ \ \  \\ MATH500} \vspace{-0.6mm} \\
    
    % \bottomrule
    \end{tabular}
}
\end{table*}
